\documentclass[11pt]{article}
\usepackage[a4paper,margin=2.35cm]{geometry}
\usepackage{fontspec}
\usepackage{amsmath,amssymb}
\setmainfont[ItalicFont=Noto Serif Italic]{Malgun Gothic}
\setsansfont{Malgun Gothic}
\setlength{\parindent}{1.2em}
\setlength{\parskip}{0.45em}
\setlength{\emergencystretch}{3em}
\newcommand{\mat}[1]{\begin{pmatrix}#1\end{pmatrix}}
\begin{document}

\emph{a)} $\Omega^*=T^{-1}\Omega T$가 체 $K$ 위의 유리 표현이라고 하자. 또 $I^*=T^{-1}IT$, $J^*=T^{-1}JT$라 하자. $J$와 $J^*$는 모두 $K$ 위에서 유리적인 서로 상사인 행렬이므로, $J^*$를 $J$로 옮기는 $K$ 위의 유리 변환 $R$도 존재한다.
\[
        R^{-1}J^*R=J.
\]
처음부터 $\Omega^*$를 역시 $K$ 위에서 유리적인 표현 $R^{-1}\Omega^*R$로 바꾸면, 일반성을 잃지 않고 $J^*=J$이라 가정할 수 있음을 알 수 있다. 이제
\[
        I^*=\mat{a&b\\ c&d}\qquad (a,b,c,d\text{는 }K\text{의 원소})
\]
라 하자. $IJ=-JI$에서 $I^*J^*=-J^*I^*$가 따르고, $J^*=J$이므로
\[
        a=-d,
        \qquad b=c
\]
를 얻는다. 또 $I^2=-E$에서 $I^{*2}=-E$가 따르므로
\[
        \mat{-1&0\\0&-1}=\mat{a&b\\ b&-a}^2
        =\mat{a^2+b^2&0\\0&a^2+b^2}.
\]
따라서 $a^2+b^2=-1$이고, 실제로 $-1$을 $K$에서 두 제곱의 합으로 나타낼 수 있다.

\end{document}
